\documentclass[a4paper,12pt]{article}

\usepackage[T2A]{fontenc}
\usepackage[utf8]{inputenc}
\usepackage[russian]{babel}
\usepackage{amsmath,amssymb,amsthm,mathtools}
\usepackage{helvet}
\renewcommand{\familydefault}{\sfdefault}
\usepackage{icomma}
\usepackage[a4paper,margin=2.1cm,headheight=25pt]{geometry}
\usepackage{microtype}
\usepackage{tikz}
\usepackage{tkz-euclide}
\usepackage{pgfplots}
\pgfplotsset{compat=1.18}
\usepackage{tabularx,booktabs,longtable}
\usepackage{enumitem}
\usepackage[hidelinks]{hyperref}
\usepackage{parskip}
\usepackage{fancyhdr}
\usepackage{setspace}
\usepackage{xcolor}

\definecolor{accent}{HTML}{008080}
\definecolor{darkaccent}{HTML}{004D4D}
\definecolor{textgray}{HTML}{333333}
\definecolor{softgray}{HTML}{F1F4F4}

\onehalfspacing
\setlength{\parindent}{0pt}
\setlength{\parskip}{0.55em}
\setlist[itemize]{leftmargin=1.45em,itemsep=0.25em,topsep=0.25em}
\setlist[enumerate]{leftmargin=1.7em,itemsep=0.55em,topsep=0.35em}
\emergencystretch=2em
\raggedbottom

\tkzSetUpLine[line width=1pt,color=accent]
\tkzSetUpPoint[color=accent,fill=accent,size=3]
\tkzSetUpLabel[color=black,font=\small]

\pagestyle{fancy}
\fancyhf{}
\fancyhead[L]{\small\color{textgray}Геометрия: медиана в треугольнике}
\fancyhead[R]{\small\color{textgray}София Кальней}
\fancyfoot[C]{\small\thepage}
\renewcommand{\headrulewidth}{0.35pt}
\renewcommand{\headrule}{%
  \hbox to\headwidth{%
    \color{accent!55}\leaders\hrule height \headrulewidth\hfill
  }%
}

\newcommand{\topic}[1]{%
  \vspace{0.35em}%
  {\large\bfseries\color{darkaccent}#1\par}%
  \vspace{0.15em}%
  {\color{accent!70}\rule{\linewidth}{0.45pt}}%
  \vspace{0.35em}%
}

\newcommand{\subtopic}[1]{%
  \vspace{0.2em}%
  {\bfseries\color{darkaccent}#1}\par
}

\newcommand{\key}[1]{\textbf{\color{darkaccent}#1}}

\newcommand{\formula}[1]{%
  \[
    \boxed{\displaystyle #1}
  \]
}

\begin{document}

% ============================================================
% ТИТУЛЬНЫЙ ЛИСТ
% ============================================================

\begin{titlepage}
\thispagestyle{empty}
\centering

\vspace*{2.2cm}

{\color{darkaccent}\rule{0.22\linewidth}{1.1pt}}\par

\vspace{0.8cm}

{\Huge\bfseries Чек-лист\par}

\vspace{0.55cm}

{\LARGE\bfseries\color{darkaccent}
Медиана в треугольнике\par}

\vspace{0.2cm}

{\Large
Свойства, теорема Аполлония и задачи на площадь\par}

\vspace{1.4cm}

{\large
София Кальней \textbullet{} ОГЭ\par}

\vfill

{\large
Лёвин Артём Александрович\par}

\vspace{0.15cm}

{\normalsize
эксклюзивно для Софии Кальней\par}

\vspace{0.55cm}

{\normalsize
\today\par}

\vfill

\begin{tikzpicture}[x=1cm,y=1cm]
  \draw[
    accent!55,
    line width=1.15pt
  ]
  (0,0)
    .. controls (2.3,0.55) and (5.0,-0.45) .. (7.1,0.05)
    .. controls (9.2,0.55) and (11.7,-0.35) .. (14.2,0.10);
\end{tikzpicture}

\end{titlepage}

% ============================================================
% ОСНОВНОЙ ЧЕК-ЛИСТ
% ============================================================

\topic{1. Медиана: что нужно узнавать сразу}

\key{Определение.}
Медиана треугольника соединяет вершину с серединой
противоположной стороны.

Если $AM$ --- медиана треугольника $ABC$, то

\[
BM=MC=\frac{BC}{2}.
\]

\key{Равные площади.}
Треугольники $ABM$ и $ACM$ имеют общую высоту,
проведённую из $A$ к прямой $BC$, и равные основания
$BM$ и $MC$.

Поэтому

\formula{
S_{ABM}=S_{ACM}=\frac12 S_{ABC}.
}

\subtopic{Быстрая проверка условия}

Если в условии фигурируют медиана и равенство вида
$AM=BM$, сразу используйте факт $BM=MC$.

Часто после этого возникают три равных отрезка,
что позволяет увидеть окружность с центром в точке $M$.

\begin{center}
\begin{tikzpicture}[scale=0.82]

  \tkzDefPoint(0,0){B}
  \tkzDefPoint(6,0){C}
  \tkzDefPoint(2.1,3.2){A}
  \tkzDefPoint(3,0){M}

  \tkzDrawPolygon[
    line width=0.9pt,
    color=textgray
  ](A,B,C)

  \tkzDrawSegment[
    line width=1.15pt,
    color=accent
  ](A,M)

  \tkzDrawPoints(A,B,C,M)

  \tkzLabelPoints[above](A)
  \tkzLabelPoints[below](B,M,C)

\end{tikzpicture}
\end{center}

% ------------------------------------------------------------

\topic{2. Медиана к гипотенузе прямоугольного треугольника}

Пусть $\angle ABC=90^\circ$, а $BM$ --- медиана
к гипотенузе $AC$.

Тогда середина гипотенузы равноудалена от всех трёх
вершин:

\formula{
BM=AM=CM=\frac{AC}{2}.
}

Верно и обратное утверждение.

Если $BM$ --- медиана к стороне $AC$ и

\[
BM=AM=CM,
\]

то точки $A$, $B$, $C$ лежат на окружности
с центром $M$, а $AC$ является диаметром.

По теореме Фалеса

\[
\angle ABC=90^\circ.
\]

\subtopic{Связь с высотой к гипотенузе}

Если катеты прямоугольного треугольника равны
$a$ и $b$, гипотенуза равна $c$, а высота
к гипотенузе равна $h$, то

\[
S=\frac12 ab
\]

и одновременно

\[
S=\frac12 ch.
\]

Следовательно,

\formula{
h=\frac{ab}{c}.
}

\subtopic{Пример: $AB=6$, $BC=8$, $BM$ --- медиана, $BK$ --- высота}

Из условия $AM=BM$ и свойства медианы имеем

\[
AM=BM=CM.
\]

Следовательно, треугольник $ABC$ прямоугольный
при вершине $B$.

По теореме Пифагора

\[
AC=\sqrt{6^2+8^2}=10.
\]

Тогда

\[
BM=\frac{AC}{2}=5.
\]

Высота к гипотенузе:

\[
BK=\frac{AB\cdot BC}{AC}
   =\frac{6\cdot8}{10}
   =\frac{24}{5}.
\]

Треугольник $BKM$ прямоугольный при $K$,
а его гипотенуза --- $BM$.

Поэтому

\formula{
\cos\angle KBM
=
\frac{BK}{BM}
=
\frac{24}{25}
=
0{,}96.
}

\begin{center}
\begin{tikzpicture}[scale=0.95]

  \tkzDefPoint(0,0){A}
  \tkzDefPoint(5,0){C}
  \tkzDefPoint(1.8,2.4){B}
  \tkzDefPoint(2.5,0){M}
  \tkzDefPoint(1.8,0){K}

  \tkzDrawPolygon[
    line width=0.9pt,
    color=textgray
  ](A,B,C)

  \tkzDrawSegment[
    line width=1.15pt,
    color=accent
  ](B,M)

  \tkzDrawSegment[
    line width=1pt,
    color=darkaccent
  ](B,K)

  \tkzMarkRightAngle[
    size=0.25,
    draw=darkaccent
  ](A,B,C)

  \tkzMarkRightAngle[
    size=0.22,
    draw=darkaccent
  ](B,K,C)

  \tkzDrawPoints(A,B,C,M,K)

  \tkzLabelPoints[below](A,K,M,C)
  \tkzLabelPoints[above](B)

\end{tikzpicture}
\end{center}

% ------------------------------------------------------------

\topic{3. Продление медианы: превращаем треугольник в параллелограмм}

Пусть $AM$ --- медиана треугольника $ABC$.

Продлим $AM$ за точку $M$ до точки $F$ так,
чтобы

\[
AM=MF.
\]

Поскольку

\[
BM=MC
\]

и

\[
AM=MF,
\]

диагонали $AF$ и $BC$ четырёхугольника $ABFC$
делятся точкой $M$ пополам.

Следовательно, $ABFC$ --- параллелограмм.

Отсюда

\[
AF=2AM,
\qquad
FC=AB,
\qquad
BF=AC.
\]

Каждая диагональ параллелограмма делит его
на два равновеликих треугольника.

Поэтому

\formula{
S_{ABC}=S_{AFC}.
}

\begin{center}
\begin{tikzpicture}[scale=0.82]

  \tkzDefPoint(0,0){A}
  \tkzDefPoint(1,2.5){B}
  \tkzDefPoint(5,0){C}
  \tkzDefPoint(6,2.5){F}
  \tkzDefPoint(3,1.25){M}

  \tkzDrawPolygon[
    line width=0.9pt,
    color=textgray
  ](A,B,F,C)

  \tkzDrawSegment[
    line width=1.15pt,
    color=accent
  ](A,F)

  \tkzDrawSegment[
    line width=1pt,
    color=darkaccent
  ](B,C)

  \tkzDrawPoints(A,B,C,F,M)

  \tkzLabelPoints[below](A,C)
  \tkzLabelPoints[above](B,F)
  \tkzLabelPoints[right](M)

\end{tikzpicture}
\end{center}

\key{Зачем нужна конструкция.}

Если известны две стороны

\[
AB=c,
\qquad
AC=b
\]

и медиана

\[
AM=m,
\]

то в треугольнике $AFC$ известны все три стороны:

\[
AC=b,
\qquad
FC=c,
\qquad
AF=2m.
\]

Значит, его площадь можно найти по формуле Герона.
Полученная площадь совпадает с площадью исходного
треугольника.

% ------------------------------------------------------------

\topic{4. Теорема Аполлония для медианы}

Пусть

\[
a=BC,
\qquad
b=CA,
\qquad
c=AB,
\]

а

\[
m_a=AM
\]

--- медиана к стороне $a$.

Для параллелограмма $ABFC$ выполняется тождество

\[
AF^2+BC^2
=
2\bigl(AB^2+AC^2\bigr).
\]

Поскольку

\[
AF=2m_a,
\]

получаем

\[
4m_a^2+a^2=2b^2+2c^2.
\]

Отсюда

\formula{
m_a=
\frac12
\sqrt{2b^2+2c^2-a^2}.
}

Если требуется найти сторону, удобно записать формулу
в виде

\formula{
a^2=2b^2+2c^2-4m_a^2.
}

\key{Контроль здравого смысла.}

Подкоренное выражение должно быть положительным,
а найденные стороны должны удовлетворять
неравенствам треугольника.

% ------------------------------------------------------------

\topic{5. Площадь по двум сторонам и медиане}

Если известны стороны $b$, $c$ и медиана $m$,
проведённая к третьей стороне, используйте
конструкцию с продлением медианы.

Получается треугольник со сторонами

\[
b,
\qquad
c,
\qquad
2m.
\]

Его полупериметр

\[
p=\frac{b+c+2m}{2}.
\]

По формуле Герона

\formula{
S=
\sqrt{
p(p-b)(p-c)(p-2m)
}.
}

Для существования такой конфигурации необходимо

\[
|b-c|<2m<b+c.
\]

\subtopic{Пример: стороны $27$ и $29$, медиана к третьей стороне равна $26$}

Продлеваем медиану вдвое.

Получаем треугольник со сторонами

\[
27,
\qquad
29,
\qquad
52.
\]

Его полупериметр

\[
p=
\frac{27+29+52}{2}
=
54.
\]

Тогда

\begin{align*}
S
&=
\sqrt{
54(54-27)(54-29)(54-52)
}
\\
&=
\sqrt{
54\cdot27\cdot25\cdot2
}
\\
&=
270.
\end{align*}

Следовательно,

\formula{
S_{ABC}=270.
}

Для проверки можно применить теорему Аполлония:

\begin{align*}
BC^2
&=
2\cdot27^2
+
2\cdot29^2
-
4\cdot26^2
\\
&=
436.
\end{align*}

Отсюда

\[
BC=2\sqrt{109}.
\]

% ------------------------------------------------------------

\topic{6. Квадраты на катетах: как увидеть скрытую высоту}

Пусть $ABC$ --- прямоугольный треугольник,

\[
\angle ACB=90^\circ.
\]

На катетах $AC$ и $BC$ снаружи построены квадраты
$ACDE$ и $BCFG$.

Медиана $CM$, проведённая к гипотенузе $AB$,
продолжена до пересечения с $DF$ в точке $N$.

Обозначим

\[
AC=a,
\qquad
BC=b,
\qquad
AB=c.
\]

По теореме Пифагора

\[
c=\sqrt{a^2+b^2}.
\]

Из свойств квадратов

\[
DC=a,
\qquad
CF=b,
\]

а также

\[
DC\parallel BC,
\qquad
CF\parallel AC.
\]

Поскольку $M$ --- середина гипотенузы,

\[
CM=AM=BM.
\]

Рассмотрим углы.

Пусть

\[
\angle ACM=\alpha.
\]

Тогда в равнобедренном треугольнике $ACM$

\[
\angle CAM=\alpha.
\]

Так как

\[
\angle ACB=90^\circ,
\]

имеем

\[
\angle MCB=90^\circ-\alpha.
\]

Из параллельности сторон квадратов соответствующие
острые углы в треугольнике $DCF$ совпадают
с острыми углами исходного прямоугольного
треугольника.

Поэтому

\[
\angle DNC=90^\circ.
\]

Следовательно,

\[
CN\perp DF.
\]

Значит, в прямоугольном треугольнике $DCF$
отрезок $CN$ является высотой к гипотенузе $DF$.

При этом

\[
DF
=
\sqrt{DC^2+CF^2}
=
\sqrt{a^2+b^2}
=
c.
\]

Из двух формул площади треугольника $DCF$:

\[
S_{DCF}
=
\frac12 ab
\]

и

\[
S_{DCF}
=
\frac12\,DF\cdot CN
\]

получаем

\formula{
CN=\frac{ab}{c}.
}

\begin{center}
\begin{tikzpicture}[scale=0.68]

  \tkzDefPoint(0,0){C}
  \tkzDefPoint(0,3){A}
  \tkzDefPoint(4,0){B}

  \tkzDefPoint(-3,0){D}
  \tkzDefPoint(-3,3){E}

  \tkzDefPoint(0,-4){F}
  \tkzDefPoint(4,-4){G}

  \tkzDefPoint(2,1.5){M}
  \tkzDefPoint(-1.92,-1.44){N}

  \tkzDrawPolygon[
    line width=0.9pt,
    color=textgray
  ](A,C,B)

  \tkzDrawPolygon[
    line width=0.8pt,
    color=accent
  ](A,C,D,E)

  \tkzDrawPolygon[
    line width=0.8pt,
    color=accent
  ](B,C,F,G)

  \tkzDrawSegment[
    line width=1pt,
    color=darkaccent
  ](D,F)

  \tkzDrawSegment[
    line width=1.15pt,
    color=accent
  ](N,M)

  \tkzMarkRightAngle[
    size=0.22,
    draw=darkaccent
  ](A,C,B)

  \tkzMarkRightAngle[
    size=0.18,
    draw=darkaccent
  ](D,N,C)

  \tkzDrawPoints(A,B,C,D,E,F,G,M,N)

  \tkzLabelPoints[above](A,E,M)
  \tkzLabelPoints[right](B,G)
  \tkzLabelPoints[below](D,N,F)
  \tkzLabelPoints[above right](C)

\end{tikzpicture}
\end{center}

Например, если

\[
a=5,
\qquad
b=12,
\]

то

\[
c=13
\]

и

\[
CN
=
\frac{5\cdot12}{13}
=
\frac{60}{13}.
\]

% ------------------------------------------------------------

\topic{7. Типичные ошибки и контроль решения}

\begin{itemize}

  \item
  Медиана в общем треугольнике не обязана быть
  высотой или биссектрисой.

  \item
  Равенство медианы половине противоположной стороны
  даёт прямой угол только при корректно указанной
  медиане к этой стороне.

  \item
  В прямоугольном треугольнике сначала определите
  гипотенузу: она лежит напротив прямого угла.

  \item
  В формуле Аполлония коэффициент при $m_a^2$
  равен $4$:

  \[
  a^2+4m_a^2=2b^2+2c^2.
  \]

  \item
  При продлении медианы новая сторона равна $2m$.
  Именно эта длина используется в формуле Герона.

  \item
  В формуле Герона

  \[
  p=\frac{a+b+c}{2}
  \]

  является полупериметром.

  \item
  После нахождения неизвестной стороны проверьте
  положительность подкоренных выражений и
  неравенства треугольника.

  \item
  При вычислении косинуса в прямоугольном
  треугольнике сначала определите гипотенузу
  рассматриваемого малого треугольника.

\end{itemize}

% ============================================================
% ТРЕНИРОВОЧНЫЙ БЛОК
% ============================================================

\clearpage

\begin{center}

{\LARGE\bfseries\color{darkaccent}
Тренировочный блок\par}

\vspace{0.15cm}

{\normalsize
10 задач по нарастающей сложности\par}

\end{center}

\vspace{0.35cm}

\subtopic{Средний уровень}

\begin{enumerate}[label=\textbf{\arabic*.}]

  \item
  В треугольнике $ABC$ отрезок $AM$ --- медиана
  к стороне $BC$.

  Известно, что

  \[
  BC=18.
  \]

  Найдите $BM$.

  \item
  В прямоугольном треугольнике $ABC$
  угол $B$ прямой,

  \[
  AB=9,
  \qquad
  BC=12.
  \]

  Медиана $BM$ проведена к гипотенузе $AC$.
  Найдите $BM$.

  \item
  В прямоугольном треугольнике $ABC$
  угол $B$ прямой,

  \[
  AB=5,
  \qquad
  BC=12.
  \]

  $BM$ --- медиана к $AC$,
  $BK$ --- высота к $AC$.

  Найдите

  \[
  \cos\angle KBM.
  \]

\end{enumerate}

\subtopic{Выше среднего}

\begin{enumerate}[label=\textbf{\arabic*.},start=4]

  \item
  В треугольнике $ABC$

  \[
  AB=13,
  \qquad
  AC=15,
  \qquad
  BC=14.
  \]

  Найдите медиану $AM$, проведённую к стороне $BC$.

  \item
  В треугольнике $ABC$

  \[
  AB=17,
  \qquad
  AC=10,
  \]

  а медиана $AM$ к стороне $BC$ равна $11$.

  Найдите $BC$.

  \item
  Две стороны треугольника равны $13$ и $15$,
  а медиана, проведённая к третьей стороне,
  равна $7$.

  Найдите площадь треугольника.

  \item
  В прямоугольном треугольнике $ABC$
  угол $C$ прямой,

  \[
  AC=8,
  \qquad
  BC=15.
  \]

  На катетах снаружи построены квадраты
  $ACDE$ и $BCFG$.

  Точка $M$ --- середина $AB$,
  а прямая $CM$ пересекает $DF$ в точке $N$.

  Найдите $CN$.

  \item
  В треугольнике $ABC$

  \[
  AB=10,
  \qquad
  AC=24,
  \]

  а медиана $AM$ к стороне $BC$ равна $13$.

  Определите вид треугольника по углам
  и найдите его площадь.

  \item
  Две стороны треугольника равны $13$ и $15$,
  а медиана к третьей стороне равна $10$.

  Найдите площадь треугольника.

\end{enumerate}

\subtopic{Сложная задача}

\begin{enumerate}[label=\textbf{\arabic*.},start=10]

  \item
  Две стороны треугольника равны $25$ и $39$,
  а медиана, проведённая к третьей стороне,
  равна $28$.

  Найдите третью сторону и площадь треугольника.

\end{enumerate}

% ============================================================
% ОТВЕТЫ
% ============================================================

\clearpage

\begin{center}

{\LARGE\bfseries\color{darkaccent}
Ответы\par}

\vspace{0.2cm}

{\small
Используйте таблицу после самостоятельной попытки.\par}

\end{center}

\vspace{0.5cm}

\renewcommand{\arraystretch}{1.4}

\begin{table}[h!]
\centering

\caption{Ответы к тренировочному блоку}

\begin{tabularx}{0.78\linewidth}{@{}cX@{}}

\toprule

\textbf{№} & \textbf{Ответ} \\

\midrule

1 & $9$ \\

2 & $\dfrac{15}{2}$ \\

3 & $\dfrac{120}{169}$ \\

4 & $2\sqrt{37}$ \\

5 & $7\sqrt{6}$ \\

6 & $84$ \\

7 & $\dfrac{120}{17}$ \\

8 & $\angle A=90^\circ$; $S=120$ \\

9 & $12\sqrt{66}$ \\

10 & третья сторона $34$; $S=420$ \\

\bottomrule

\end{tabularx}
\end{table}

\vfill

{\small\color{textgray}
Финальная самопроверка:
определили, к какой стороне проведена медиана;
выбрали нужное свойство;
проверили длины, радикалы и неравенства треугольника;
записали ответ в требуемом виде.
}

\end{document}